\chapter{Linear Independence, Span, Bases, and Dimension}
\label{mf:ch06-linear-independence-span-bases-dimension}

The previous chapters treated vectors as ordered collections of quantities and described how to add, subtract, scale, combine, compare, and measure them. We now ask a structural question: when does a collection of vectors contain genuinely different directions, and when is one vector redundant because it can be built from the others?

This chapter develops four connected ideas. The \emph{span} of a set of vectors is the collection of all vectors that can be built from them by linear combinations. \emph{Linear independence} tells us whether any vector in the collection is redundant. A \emph{basis} is a nonredundant set of directions that can represent every vector in a space. The \emph{dimension} of a space counts how many independent directions are needed to describe it.

These ideas will later support matrix rank, invertibility, least squares, degrees of freedom, and identifiability. For now, we stay with vectors and linear combinations.

\section{Span: All Vectors Built from a Collection}
\label{mf:span}

Let \(\mathbf{a}_1,\ldots,\mathbf{a}_k\in\mathbb{R}^n\). The \emph{span} of these vectors is the set of all linear combinations of them:
\[
    \operatorname{span}\{\mathbf{a}_1,\ldots,\mathbf{a}_k\}
    =
    \left\{\beta_1\mathbf{a}_1+\cdots+\beta_k\mathbf{a}_k:
    \beta_1,\ldots,\beta_k\in\mathbb{R}\right\}.
\]
In words, the span is everything we can make by scaling and adding the listed vectors.

If \(\mathbf{a}_1\) is a single nonzero vector, then
\[
    \operatorname{span}\{\mathbf{a}_1\}=\{\beta\mathbf{a}_1:\beta\in\mathbb{R}\},
\]
which is the line through the origin in the direction of \(\mathbf{a}_1\). If two nonzero vectors in \(\mathbb{R}^3\) point in genuinely different directions, their span is a plane through the origin. If three vectors in \(\mathbb{R}^3\) point in enough different directions, their span may be all of \(\mathbb{R}^3\).

The span language makes linear combinations geometric. A set of vectors supplies directions; the span is the collection of all locations reachable by moving in those directions with real-valued coefficients.

\section{Linear Dependence}
\label{mf:linear-dependence}

A set of vectors \(\{\mathbf{a}_1,\ldots,\mathbf{a}_k\}\subset\mathbb{R}^n\) is \emph{linearly dependent} if there exist scalars \(\beta_1,\ldots,\beta_k\), not all zero, such that
\[
    \beta_1\mathbf{a}_1+\beta_2\mathbf{a}_2+\cdots+\beta_k\mathbf{a}_k=\mathbf{0}.
\]
The phrase ``not all zero'' is essential. We can always make the zero vector by choosing all coefficients equal to zero. Linear dependence means there is a \emph{nontrivial} way to make the zero vector.

For two vectors, suppose
\[
    \beta_1\mathbf{a}_1+\beta_2\mathbf{a}_2=\mathbf{0}
\]
with \(\beta_2\neq 0\). Then
\[
    \mathbf{a}_2=-\frac{\beta_1}{\beta_2}\mathbf{a}_1.
\]
Thus \(\mathbf{a}_2\) is just a scalar multiple of \(\mathbf{a}_1\). The two vectors do not provide two different directions; one is redundant.

The same idea holds for more than two vectors. If a collection is linearly dependent, then at least one vector in the collection can be written as a linear combination of the others. To see why, suppose
\[
    \beta_1\mathbf{a}_1+\cdots+\beta_k\mathbf{a}_k=\mathbf{0}
\]
for coefficients not all zero. Choose an index \(j\) with \(\beta_j\neq 0\). Then
\[
    \mathbf{a}_j
    =
    -\frac{1}{\beta_j}\sum_{i\neq j}\beta_i\mathbf{a}_i.
\]
So \(\mathbf{a}_j\) can be built from the remaining vectors.

A collection that contains the zero vector is automatically linearly dependent. If \(\mathbf{a}_j=\mathbf{0}\), choose \(\beta_j=1\) and all other coefficients equal to zero. That is a nontrivial linear combination that equals \(\mathbf{0}\).

\section{Linear Independence}
\label{mf:linear-independence}

A set of vectors \(\mathbf{a}_1,\ldots,\mathbf{a}_k\in\mathbb{R}^n\) is \emph{linearly independent} if it is not linearly dependent. Equivalently,
\[
    \beta_1\mathbf{a}_1+\cdots+\beta_k\mathbf{a}_k=\mathbf{0}
    \quad\Longrightarrow\quad
    \beta_1=\beta_2=\cdots=\beta_k=0.
\]
In words, the only way to combine the vectors to get the zero vector is the trivial way: use all-zero coefficients.

Linear independence means that no vector in the set is unnecessary. Each vector contributes a direction that cannot be built from the others.

\paragraph{Example in \(\mathbb{R}^2\).}
Consider
\[
    \mathbf{a}_1=\begin{bmatrix}1\\0\end{bmatrix},
    \qquad
    \mathbf{a}_2=\begin{bmatrix}1\\1\end{bmatrix}.
\]
If
\[
    \beta_1\mathbf{a}_1+\beta_2\mathbf{a}_2=\mathbf{0},
\]
then
\[
    \beta_1\begin{bmatrix}1\\0\end{bmatrix}
    +\beta_2\begin{bmatrix}1\\1\end{bmatrix}
    =
    \begin{bmatrix}\beta_1+\beta_2\\ \beta_2\end{bmatrix}
    =
    \begin{bmatrix}0\\0\end{bmatrix}.
\]
The second coordinate gives \(\beta_2=0\). Then the first coordinate gives \(\beta_1=0\). Therefore the two vectors are linearly independent.

\section{How Many Independent Vectors Can There Be?}
\label{mf:maximal-independent-sets}
\label{mf:dimension}

A linearly independent set of \(n\)-vectors can have at most \(n\) vectors. In \(\mathbb{R}^2\), we cannot have three linearly independent vectors. In \(\mathbb{R}^3\), we cannot have four linearly independent vectors. More generally, no more than \(n\) independent directions fit inside \(\mathbb{R}^n\).

One way to see this is to examine the equation
\[
    \beta_1\mathbf{a}_1+\cdots+\beta_k\mathbf{a}_k=\mathbf{0}.
\]
This is a system of \(n\) scalar equations, one for each coordinate, in \(k\) unknown coefficients \(\beta_1,\ldots,\beta_k\). If \(k>n\), then there are more unknown coefficients than equations. Such a homogeneous system has a nontrivial solution, so the vectors are linearly dependent.

This fact gives the first meaning of \emph{dimension}: \(\mathbb{R}^n\) has \(n\) independent directions. The number \(n\) is not just the number of entries in a vector; it is the maximum number of mutually independent directions available in the space.

\section{Bases}
\label{mf:bases}

A \emph{basis} for \(\mathbb{R}^n\) is a collection of \(n\) linearly independent vectors in \(\mathbb{R}^n\). Equivalently, it is a set of vectors that is both:
\begin{enumerate}
    \item linearly independent, so none of the vectors is redundant; and
    \item spanning, so every vector in \(\mathbb{R}^n\) can be written as a linear combination of them.
\end{enumerate}

Suppose \(\mathbf{a}_1,\ldots,\mathbf{a}_n\) form a basis for \(\mathbb{R}^n\). Then every vector \(\mathbf{b}\in\mathbb{R}^n\) can be written in the form
\[
    \mathbf{b}=c_1\mathbf{a}_1+c_2\mathbf{a}_2+\cdots+c_n\mathbf{a}_n
\]
for some scalars \(c_1,\ldots,c_n\).

The source notes justify this using the maximum-size fact. The set
\[
    \{\mathbf{a}_1,\ldots,\mathbf{a}_n,\mathbf{b}\}
\]
contains \(n+1\) vectors in \(\mathbb{R}^n\), so it must be linearly dependent. Therefore there are coefficients, not all zero, such that
\[
    \beta_1\mathbf{a}_1+\cdots+\beta_n\mathbf{a}_n+\beta_{n+1}\mathbf{b}=\mathbf{0}.
\]
The coefficient \(\beta_{n+1}\) cannot be zero; if it were zero, then the basis vectors themselves would be linearly dependent. Thus \(\beta_{n+1}\neq0\), and we may solve for \(\mathbf{b}\):
\[
    \mathbf{b}
    =
    -\frac{\beta_1}{\beta_{n+1}}\mathbf{a}_1
    -\cdots-
    \frac{\beta_n}{\beta_{n+1}}\mathbf{a}_n.
\]
So every vector in \(\mathbb{R}^n\) is in the span of the basis vectors.

\section{Coordinates in a Basis Are Unique}
\label{mf:coordinate-representation}

A basis does more than merely represent every vector. It represents every vector \emph{uniquely}.

Suppose \(\mathbf{a}_1,\ldots,\mathbf{a}_n\) form a basis and suppose that \(\mathbf{b}\in\mathbb{R}^n\) has two representations:
\[
    \mathbf{b}=c_1\mathbf{a}_1+\cdots+c_n\mathbf{a}_n
\]
and
\[
    \mathbf{b}=d_1\mathbf{a}_1+\cdots+d_n\mathbf{a}_n.
\]
Subtract the second equation from the first:
\[
    \mathbf{0}
    =
    (c_1-d_1)\mathbf{a}_1+\cdots+(c_n-d_n)\mathbf{a}_n.
\]
Because the basis vectors are linearly independent,
\[
    c_1-d_1=\cdots=c_n-d_n=0.
\]
Therefore \(c_i=d_i\) for every \(i\). The coordinates of \(\mathbf{b}\) in that basis are unique.

This is the reason bases are useful. Once a basis has been chosen, a vector can be described by its coordinate weights relative to that basis.

\section{The Standard Basis}
\label{mf:standard-basis}

The standard basis for \(\mathbb{R}^n\) is
\[
    \mathbf{e}_1=\begin{bmatrix}1\\0\\\vdots\\0\end{bmatrix},
    \quad
    \mathbf{e}_2=\begin{bmatrix}0\\1\\\vdots\\0\end{bmatrix},
    \quad \ldots \quad,
    \mathbf{e}_n=\begin{bmatrix}0\\0\\\vdots\\1\end{bmatrix}.
\]
Each \(\mathbf{e}_i\) has a 1 in the \(i\)th entry and 0 everywhere else. For any vector
\[
    \mathbf{b}=\begin{bmatrix}b_1\\b_2\\\vdots\\b_n\end{bmatrix}\in\mathbb{R}^n,
\]
we have
\[
    \mathbf{b}=b_1\mathbf{e}_1+b_2\mathbf{e}_2+\cdots+b_n\mathbf{e}_n.
\]
Thus the standard basis expresses the usual coordinates of \(\mathbf{b}\).

Other bases are possible. For example, the two vectors
\[
    \begin{bmatrix}1.2\\-2.6\end{bmatrix},
    \qquad
    \begin{bmatrix}-0.3\\-3.7\end{bmatrix}
\]
form a basis for \(\mathbb{R}^2\). They point in different directions, so neither one is a scalar multiple of the other. Therefore they provide two independent directions in a two-dimensional space.

\section{Orthonormal Bases}
\label{mf:orthonormal-basis-note}

The standard basis is special because its vectors are both orthogonal and normal. For the standard basis vectors,
\[
    \mathbf{e}_i^T\mathbf{e}_j=
    \begin{cases}
        1, & i=j,\\
        0, & i\neq j.
    \end{cases}
\]
A basis whose vectors are mutually orthogonal and each have norm 1 is called an \emph{orthonormal basis}. Equivalently, for an orthonormal basis \(\mathbf{q}_1,\ldots,\mathbf{q}_n\),
\[
    \mathbf{q}_i^T\mathbf{q}_j=
    \begin{cases}
        1, & i=j,\\
        0, & i\neq j.
    \end{cases}
\]

For an orthonormal basis, dot products recover coordinates directly. These bases return later with Gram--Schmidt and QR factorization. This chapter only introduces the idea; the construction is deferred.

\section{Data Analysis Bridge}
\label{mf:ch06-data-analysis-bridge}

The Data Analysis textbook uses the ideas in this chapter in several places:
\begin{itemize}
    \item A model becomes more flexible when it has more independent directions available for fitting patterns.
    \item Degrees of freedom count independent dimensions, not just the number of formulas on a page.
    \item Regression design columns must provide enough independent information for coefficients to be uniquely identified.
    \item Multicollinearity is a practical warning that some design columns are nearly, or exactly, redundant.
\end{itemize}

Later chapters reuse this vocabulary for matrix rank, column space, null space, and least-squares identifiability. Here the key point is simpler: independent vectors add directions, dependent vectors repeat directions, bases give unique coordinates, and dimension counts directions.

\section{Chapter Summary}
\label{mf:linear-independence-summary}

Carry forward these structure ideas:
\begin{itemize}
    \item The span of a set of vectors is the collection of all linear combinations of those vectors.
    \item A linearly dependent set contains redundancy; a linearly independent set does not.
    \item No more than \(n\) linearly independent vectors can exist in \(\mathbb{R}^n\).
    \item A basis for \(\mathbb{R}^n\) is an independent set that represents every vector in \(\mathbb{R}^n\), and those coordinates are unique.
    \item The standard basis is the default coordinate system; orthonormal bases will become useful when we build Gram-Schmidt and QR.
\end{itemize}
